\section{Finite valuation extensions and the actual normalized arithmetic choice}
\label{sec:arithmetic-choice}

This section constructs the arithmetic comparison required by
\texttt{NChoice}, and then its exact character and tuple ties.
Throughout, \(D\) is a supplied \texttt{D7}. Write
\[
 L=D.\texttt{what}_2.\texttt{LL},\qquad L_0=D.\texttt{LL0},
 \qquad m=D.R.m,\qquad
 W_{\rm tw}=\{w:\mu_{\rm At}(w)\ne\mu_0\}.
\]
Thus \(L\) is a number field Galois over \(\mathbb Q\), equipped
with its specified complex embedding and conjugation. Its places
are the specified prime ideals of \(\mathcal O_L\).
The additive valuation has the convention \(v_w(0)=+\infty\);
\(v_w(x)\geq0\) means integrality, and \(v_w(x)>0\) means
membership in the maximal ideal.

The rings, spaces and lattices already carried by \(D\) are
retained as supplied objects. In particular, this section does
not construct \texttt{NTw}, \texttt{NExp}, or the row model.
It proves the finite DVR extension and integral comparison
needed to assemble their \texttt{NChoice}. The equality of the
two ranks, named \texttt{RankTie}, is a separate premise.
The resulting semilinear maps preserve the specified lattice
spans; their definition makes no assertion about Hecke operators
or Hermitian forms.

\paragraph{NC1: finite extensions of a complete DVR in characteristic zero.}
Let \(A\) be a complete DVR, with uniformizer \(\pi\) and fraction
field \(K\) of characteristic zero. Let \(K'/K\) be finite and put
\(B=\{x\in K':x\text{ is integral over }A\}\).
Then \(B\) is a complete DVR, finite and free over \(A\);
its fraction field is \(K'\), and \(A\longrightarrow B\) is an
injective local homomorphism.

Here is a direct proof of the finiteness that is needed in this
application. The finite extension is separable. Choose a primitive
element, multiply it by a nonzero element of \(A\) to make it
integral, and denote the result by \(\alpha\).
The order \(C=A[\alpha]\) is free of rank \([K':K]\), with
basis \(1,\alpha,\ldots,\alpha^{d-1}\).
For \(x\in B\), the numbers
\(\operatorname{Tr}_{K'/K}(x\alpha^i)\) belong to \(A\).
Indeed, they are integral over \(A\), as sums of conjugates of
integral elements, and they belong to \(K\); the DVR \(A\) is
integrally closed. The trace pairing is nondegenerate, so
\[
 C^\vee=\{x\in K':\operatorname{Tr}_{K'/K}(xC)\subset A\}
\]
is a free \(A\)-module of rank \(d\): its coordinate matrix is
the inverse of the nonsingular trace matrix of the displayed
basis. Thus \(C\subset B\subset C^\vee\).
A submodule of a finite free module over a DVR is finite free.
It follows that \(B\) is finite over \(A\), and that
\(B[1/\pi]=K'\).

Every maximal ideal of the integral finite \(A\)-algebra \(B\)
lies over \((\pi)\). The algebra \(B/\pi B\) is finite dimensional
over \(A/\pi A\), hence Artinian. If it had two local factors,
it would have a nontrivial idempotent \(\bar e\).
Such an idempotent lifts to \(B\). To see the lifting explicitly,
start with any lift \(e_0\); when \(e_r^2-e_r\in\pi^a B\),
\((2e_r-1)^2=1+4(e_r^2-e_r)\) is a unit and
\[
 e_{r+1}=e_r-(e_r^2-e_r)(2e_r-1)^{-1}
 \quad\text{satisfies}\quad
 e_{r+1}^2-e_{r+1}\in\pi^{2a}B.
 \tag{NC.1}
\]
The inverse of \(1+\pi b\) is its convergent geometric series.
Completeness here follows from the finite free \(A\)-module
structure of \(B\). The sequence in (NC.1) converges to an
idempotent with residue \(\bar e\), contradicting that \(B\)
is a domain. Thus \(B\) is local.

For completeness, the usual DVR criterion can also be made
explicit. A nonzero prime of \(B\) contracts to a nonzero prime
of \(A\), because \(B[1/\pi]\) is a field; therefore \(B\) has
dimension one. It is Noetherian and integrally closed.
If \(\mathfrak n\) is its maximal ideal and \(0\ne a\in\mathfrak n\),
the Artinian local ring \(B/aB\) has a nonzero socle. Choose a
representative \(u\notin aB\) with
\(\mathfrak n u\subset aB\), and put \(t=u/a\).
Then \(t\notin B\) but \(t\mathfrak n\subset B\).
The inclusion \(t\mathfrak n\subset\mathfrak n\) would make \(t\)
integral over \(B\), by applying the determinant trick to a finite
generating set of the nonzero ideal \(\mathfrak n\).
Consequently \(t\mathfrak n=B\), so \(\mathfrak n=t^{-1}B\)
is principal. Since \(\mathfrak n^c\subset\pi B\) for some \(c\),
the intersection of all its powers is zero. Repeated division
by its generator expresses every nonzero element as a unit times
a unique nonnegative power. This proves the DVR assertion.
The inclusions
\(\pi B\subset\mathfrak n\) and \(\mathfrak n^c\subset\pi B\)
also show that the two adic topologies agree, so \(B\) is complete
for its maximal ideal. Its maximal ideal contracts to \((\pi)\),
which proves locality of \(A\longrightarrow B\).

This proves the characteristic-zero instance of the integral
closure supplier used in the pinned source. That instance is
enough: every coefficient fraction field below contains \(L_0\).

\paragraph{NC2: an embedding reads a place exactly when it reads its maximal ideal.}
Let \(F\) be a number field, \(B\) a DVR with fraction field \(K'\),
and \(e:F\longrightarrow K'\) a field embedding.
Every algebraic integer of \(F\) maps into \(B\): it satisfies a
monic polynomial over \(\mathbb Z\subset B\), and \(B\) is
integrally closed. This gives a unique ring map
\(\phi:\mathcal O_F\longrightarrow B\).
Suppose a nonzero prime \(\mathfrak p\subset\mathcal O_F\)
satisfies
\[
 \phi^{-1}(\mathfrak m_B)=\mathfrak p.
 \tag{NC.2}
\]
Then
\[
 v_{\mathfrak p}(x)\geq0
 \quad\Longleftrightarrow\quad e(x)\in B
 \qquad(x\in F).
 \tag{NC.3}
\]
Indeed, \(\mathcal O_{F,\mathfrak p}\) is the valuation ring at
\(\mathfrak p\). If \(x=a/s\) with \(s\notin\mathfrak p\),
then \(\phi(s)\) is a unit, giving the forward implication.
Conversely, if \(v_{\mathfrak p}(x)<0\), the element \(x^{-1}\)
belongs to its maximal ideal. Write \(x^{-1}=a/s\) with
\(a\in\mathfrak p\), \(s\notin\mathfrak p\).
Its image lies in \(\mathfrak m_B\), so its inverse cannot
belong to \(B\). This proves (NC.3), including \(x=0\).
For nonzero \(x\), the same argument gives
\(v_{\mathfrak p}(x)>0\) precisely when \(e(x)\in\mathfrak m_B\).

We also need transitivity on primes, with the direction of
conjugation specified. If \(F/F_0\) is finite Galois and
\(\mathfrak p,\mathfrak q\) lie over the same prime of
\(\mathcal O_{F_0}\), there is
\(\sigma\in\operatorname{Gal}(F/F_0)\) with
\(\sigma^{-1}\mathfrak q=\mathfrak p\).
If not, the Chinese remainder theorem gives an integer \(a\)
which is zero modulo \(\mathfrak q\) and one modulo every
conjugate of \(\mathfrak p\). Its norm
\(\prod_\tau\tau(a)\) lies in
\(\mathfrak q\cap\mathcal O_{F_0}\), and hence in
\(\mathfrak p\), although every factor avoids \(\mathfrak p\).
This contradicts primality. This argument proves precisely
the prime alignment used next; it needs no class-number premise.

\paragraph{NC3: the finite coefficient extension at a prescribed prime.}
Let \(L/\mathbb Q\) be a finite Galois extension, \(L_0\subset L\)
an intermediate field, \(A\) a complete DVR and
\(i:L_0\longrightarrow K=\operatorname{Frac}(A)\) an embedding.
Let \(\mathfrak p\) be a nonzero prime of \(\mathcal O_L\).
Assume that the maximal ideal pulled back through \(i\) on
\(\mathcal O_{L_0}\) is exactly
\(\mathfrak p\cap\mathcal O_{L_0}\).
There exist a finite complete DVR extension \(A\subset B\)
and an embedding \(e:L\longrightarrow\operatorname{Frac}(B)\)
such that
\[
 e|_{L_0}=i,\qquad
 v_{\mathfrak p}(x)\geq0\ \Longleftrightarrow\ e(x)\in B.
 \tag{NC.4}
\]
Here \(i\) in the first equality is followed by the fraction-field
embedding of \(A\subset B\). In particular, for \(x\in L_0\)
and \(y,z\in A\),
\[
 i(x)z=y\quad\Longrightarrow\quad e(x)z=y
 \quad\text{in }\operatorname{Frac}(B).
 \tag{NC.5}
\]

Choose a primitive element for \(L/L_0\). A root of its minimal
polynomial, with coefficients mapped by \(i\), in an algebraic
closure of \(K\) gives an \(L_0\)-embedding of \(L\) into a
finite extension \(K'/K\). Equivalently, adjoin the images of
a finite \(L_0\)-basis; every such image is algebraic over \(K\).
Apply NC1 to \(K'\), and call the resulting embedding \(e_1\).
The inverse image of \(\mathfrak m_B\) on \(\mathcal O_L\)
is a prime \(\mathfrak q\) above
\(\mathfrak p\cap\mathcal O_{L_0}\). It is nonzero because this
contracted prime contains a rational prime. Since \(L/L_0\)
is Galois, NC2 supplies \(\sigma\) with
\(\sigma^{-1}\mathfrak q=\mathfrak p\).
Set \(e=e_1\circ\sigma\). This changes the prime read by the
embedding and fixes its restriction to \(L_0\).
NC2 now proves (NC.4), and restriction to \(L_0\) proves (NC.5).

The literal source first adjoins a square root as well. It calls
that construction with \(a=1\) here, so this extra generator is
unnecessary; adjoining it would still give the same finite
extension argument. The assertion concerns valuation rings,
not an equality of numerical normalizations of two valuations.
A further ramification index may occur in \(B\).

\paragraph{NC4: actual readings, the residue characteristic, and the cover implication.}
For a twist index \(i\) and side \(b\in\{\mathrm{false},\mathrm{true}\}\),
put
\[
 A_0=O_0(i),\qquad A_k=O_k(i),\qquad
 s_{i,b}=\begin{cases}\mathrm{emb}_0(i)&b=\mathrm{true},\\
                       \mathrm{emb}_{0b}(i)&b=\mathrm{false}.
          \end{cases}
\]
The condition \(\texttt{ReadsAt}(i,b,w)\) says that for every
\(x\in\mathcal O_{L_0}\),
\[
 s_{i,b}(x)\in\mathfrak m_{A_0}
       \quad\Longleftrightarrow\quad v_w(x)>0.
 \tag{NC.6}
\]
The image on the left is interpreted through the unique integral
map of NC2. This is an equality of the prime being read.

The map \(A_0\to O(i)\) is local. If a nonunit had unit image,
write it as a unit times a positive power of a uniformizer;
the uniformizer would have unit image. The nonzero integer
\(\ell(i)\) is also a unit times a power of that uniformizer,
so its image would be a unit. This contradicts the supplied
condition \(\ell(i)\in\mathfrak m_{O(i)}\).
The map \(O(i)\to A_k\) is injective and finite, hence integral
and local. Moreover \(A_k\) is complete: it is finite torsion
free over the complete DVR \(O(i)\), hence a finite free
complete module; its maximal-ideal and base-adic topologies
agree because the corresponding quotient is Artinian.
Consequently
\(\phi_k:A_0\to A_k\) is injective and local.
It therefore extends to an embedding of fraction fields.

Apply NC3 to the composite
\[
 L_0\xrightarrow{s_{i,b}}\operatorname{Frac}(A_0)
       \longrightarrow\operatorname{Frac}(A_k)
\]
and the specified prime of \(L\) at \(w\).
Locality and (NC.6) verify its prime hypothesis.
This constructs, for every \(k\geq0\), a
\(\rho:\texttt{Reading}(i,k,b,w)\) consisting of a DVR \(B\),
an injective finite local \(A_k\)-algebra structure, and
\(e_\rho:L\to\operatorname{Frac}(B)\), such that
\[
 v_w(x)\geq0\Longleftrightarrow e_\rho(x)\in B,\qquad
 s_{i,b}(x)z=y\Longrightarrow
 e_\rho(x)\phi_k(z)=\phi_k(y).
 \tag{NC.7}
\]
The second equality is followed into \(B\) and its fraction field.
It holds even for \(z=0\); the later \texttt{NChTie} field
requires it only for \(z\ne0\).

Equation (NC.6) applied to \(\ell(i)\) also proves
\(\ell(i)=\operatorname{resChar}(w)\).
If two distinct rational primes belonged to the same prime
ideal, their integer Bezout combination would place \(1\)
there, which is impossible.
Finally, the supplied \texttt{TwCover} reads every prime of
\(L_0\) with side \(\mathrm{true}\). Restricting \(w\) to \(L_0\)
therefore supplies \texttt{ReadsAt} at every \(w\), and in
particular the exact off-twist condition
\(\texttt{ReadElsewhere}(D)\).
This proves the implication from the cover; it does not
construct \texttt{TwCover}.

\paragraph{NC5: the finite rational generator set of a complex form.}
Let \(V\) be an \(L\)-vector space, \(M\) a finite-dimensional
complex vector space, and \(j:V\to M\) semilinear for an
embedding \(\iota:L\to\mathbb C\).
Suppose an \(L\)-basis \((b_a)\) maps to a complex basis.
Let \(T\subset M\) be finite, span \(M\), and lie in \(j(V)\).
Then \(V\) is finite dimensional, \(j\) is injective, and
\[
 S=\{v\in V:j(v)\in T\}
 \quad\text{is finite with}\quad \operatorname{span}_L S=V.
 \tag{NC.8}
\]
The basis index is finite because its image is independent in
the finite-dimensional space \(M\).
In basis coordinates \(j\) applies \(\iota\) to every coefficient;
thus it is injective, and \(S\) has at most \(|T|\) elements.
If \(S\) failed to span, a nonzero \(L\)-linear functional
\(\lambda\) would vanish on it. Define a complex linear
functional by
\(\lambda_{\mathbb C}(j(b_a))=\iota(\lambda(b_a))\).
Semilinearity gives
\(\lambda_{\mathbb C}(j(v))=\iota(\lambda(v))\).
It vanishes on \(T\), hence on \(M\), and then on every
\(j(b_a)\), contradicting \(\lambda\ne0\).

For the literal row objects take
\[
 V=M_{t,L}(\mu,j),\quad M=M_t(\mu,j),\quad
 T=\texttt{MtLat}(\mu,j).\texttt{gen}.
 \tag{NC.9}
\]
The premises are exactly \texttt{jMt\_basis}, the finite
generator and spanning fields of the supplied lattice, and
\texttt{MtLat\_sub}. Thus the preimage of the actual displayed
generator set is finite and spans over \(L\).
This proves rational descent of those generators; it makes
no new claim that the row model itself exists.

\paragraph{NC6: scalar extension preserves the rank and the full lattice.}
Let \(A\hookrightarrow B\) be an injective map of domains,
write \(K=\operatorname{Frac}(A)\), \(K'=\operatorname{Frac}(B)\),
and let \(V\) be a \(K\)-vector space.
If \(\Lambda\subset V\) is a finitely generated \(A\)-submodule
spanning \(V\) over \(K\), put
\[
 V'=K'\otimes_A V,\qquad
 \Lambda'=\operatorname{span}_B\{1\otimes v:v\in\Lambda\}.
 \tag{NC.10}
\]
Then
\[
 \dim_{K'}V'=\dim_K V,\qquad
 \Lambda'\text{ is finite and spans }V'\text{ over }K'.
 \tag{NC.11}
\]
The injection extends to \(K\hookrightarrow K'\), and
\(K'\otimes_A V\cong K'\otimes_K V\): every nonzero element
of \(A\) already acts invertibly on both modules.
A \(K\)-basis therefore becomes a \(K'\)-basis, proving the
dimension statement. If \(\lambda_1,\ldots,\lambda_s\)
generate \(\Lambda\), then \(1\otimes\lambda_i\) generate
\(\Lambda'\) over \(B\).
For any \(v\in V\), clear denominators in a finite expression
of \(v\) in the \(\lambda_i\). There is \(0\ne d\in A\)
with \(dv\in\Lambda\), and
\[
 a\otimes v=(a/d)(1\otimes dv).
 \tag{NC.12}
\]
This proves the spanning assertion for every pure tensor
and hence for all of \(V'\).
In particular, if \(B\) is a DVR, \(\Lambda'\) is a full
finite free \(B\)-lattice.

\paragraph{NC7: matching full DVR lattices gives the required semilinear map.}
Let \(B\) be a DVR with fraction field \(K'\),
\(e:L\hookrightarrow K'\), and let \(V\) be a finite-dimensional
\(L\)-space. Suppose \(S\subset V\) is finite and spans \(V\).
Let \(U\) be a \(K'\)-space with full finite \(B\)-lattice
\(\Lambda_U\), and suppose
\(\dim_LV=\dim_{K'}U=r\).
There is an \(e\)-semilinear map \(f:V\to U\) such that
\[
 \operatorname{span}_B f(S)=\Lambda_U,
 \qquad
 \{f(b_a)\}_a\text{ is a }K'\text{-basis of }U
 \text{ for every }L\text{-basis }(b_a)\text{ of }V.
 \tag{NC.13}
\]

Choose one \(L\)-basis and let \(f_0:V\to(K')^r\) apply \(e\)
to its coordinates. The module
\(A_S=\operatorname{span}_B f_0(S)\) is finitely generated,
torsion free, and spans \((K')^r\). A finite torsion-free
module over a DVR is free; clearing denominators in a relation
shows that a \(B\)-basis is also a \(K'\)-basis. Thus both
\(A_S\) and \(\Lambda_U\) have \(B\)-bases of cardinality \(r\).
Map these bases to one another and extend \(K'\)-linearly,
obtaining \(g:(K')^r\overset{\sim}{\longrightarrow}U\).
Then \(f=g\circ f_0\) has the asserted integral span.
For any other \(L\)-basis, its change-of-basis matrix has
nonzero determinant. Applying the field embedding to that
determinant keeps it nonzero, proving the assertion for
every basis.

One elementary proof of the freeness used here is to take
a generating family of smallest cardinality. In a nontrivial
linear relation, choose a nonzero coefficient of smallest
valuation; divide all coefficients by it. Torsion freeness
allows cancellation, and the chosen coefficient becomes a
unit. One generator is then a combination of the others,
contradicting minimality. This proves freeness without an
assumed integral comparison map.
The induced map \(K'\otimes_{L,e}V\to U\) is an isomorphism;
\(f\) itself need not be surjective as a map of sets.

\paragraph{NC8: the literal \texttt{exists\_formsIso} conclusion.}
Let \(A\hookrightarrow B\), \(V\), and \(\Lambda\) be as in
NC6, with \(B\) a DVR. Let \(e:L\to\operatorname{Frac}(B)\).
For a tuple \(\mu\), index \(j\), and any proposition \(c\),
assume
\[
 c\Longrightarrow
 \dim_L M_{t,L}(\mu,j)=\dim_{\operatorname{Frac}(A)}V.
 \tag{NC.14}
\]
Then there is an \(e\)-semilinear map
\[
 f:M_{t,L}(\mu,j)\longrightarrow
       \operatorname{Frac}(B)\otimes_A V
 \tag{NC.15}
\]
which, whenever \(c\) holds, carries every \(L\)-basis to a
basis after extension of scalars and satisfies
\[
 \operatorname{span}_B
 f\bigl(\{x:j_{Mt}(x)\in\texttt{MtLat}(\mu,j).\texttt{gen}\}\bigr)
 =
 \operatorname{span}_B\{1\otimes\lambda:\lambda\in\Lambda\}.
 \tag{NC.16}
\]
If \(c\) is false, take \(f=0\). Otherwise NC5 gives the
finite spanning generator set on the left, NC6 gives the
full target lattice and its rank, and NC7 constructs \(f\).
This proves both conclusions simultaneously on the actual
generator sets. The only equality of dimensions used is
(NC.14); it is not produced by the construction.

\paragraph{NC9: all fields of the normalized \texttt{NChoice}.}
A twist packet is
\[
 p_w=(i_w,Q_w,a_w,\chi_w),\qquad
 \chi_w:\mathrm{RCG}(i_w,Q_w)\longrightarrow O_{a_w}(i_w)^\times.
\]
The source convention is
\[
 \chi^{[b]}=
 \begin{cases}\chi&b=\mathrm{true},\\
                \chi^{-1}&b=\mathrm{false},
 \end{cases}
 \qquad(\chi^{[b]})^{[b]}=\chi.
 \tag{NC.17}
\]
Suppose, for every \(w\), we have a side \(b_w\), a reading
\(\rho_w:\texttt{Reading}(i_w,a_w,b_w,w)\), and
\(\texttt{ReadsAt}(i_w,b_w,w)\).
For \(w\in W_{\rm tw}\), assume the exact rank equality
\[
 \begin{split}
 \dim_L M_{t,L}(\mu_{\rm Tw}(w),j)
 &=
 \dim_{\operatorname{Frac}(O_{a_w}(i_w))}
 M_t^{\rm loc}\bigl(i_w,j,b_w,a_w,Q_w,
                 \chi_w^{[b_w]}\circ\operatorname{art}_{i_w,Q_w}\bigr),
 \\
 &\hspace{100pt}1\leq j\leq m.
 \end{split}
 \tag{NC.18}
\]
This is exactly \(\texttt{RankTie}(p_w,b_w,w)\).

Set \(B_w=\rho_w.O_e\), \(K_w=\operatorname{Frac}(B_w)\)
and \(e_w=\rho_w.\mathrm{emb}\).
For each \(w,j\), apply NC8 to the actual local
\(\texttt{MtLat}\) and to the condition
\[
 c_{w,j}=(w\in W_{\rm tw})\ \wedge\ (1\leq j\leq m).
 \tag{NC.19}
\]
Its finite generation and spanning hypotheses are the supplied
\texttt{NExp.MtLat\_fg} and \texttt{MtLat\_span}; its rank
hypothesis is (NC.18). Choose the resulting map as
\(\texttt{eMt}(w,j)\).
Outside (NC.19) it may be the zero map, exactly as in NC8.

The full record has the following 21 fields. This list includes
the typeclass fields, whose content matters for the construction.
\[
\begin{array}{ll}
 \texttt{twI},\texttt{twI\_l},\texttt{sd}
     & i_w,\ \ell(i_w)=\operatorname{resChar}(w),\ b_w,\\
 \texttt{Q},\texttt{k'},\texttt{psi'}
     & Q_w,\ a_w,\ \chi_w^{[b_w]},\\
 \texttt{Oe},\texttt{crOe},\texttt{domOe},\texttt{dvrOe}
     & B_w\text{ and its ring, domain and DVR structures},\\
 \texttt{algOe},\texttt{algOOe},\texttt{towOe}
     & O_{a_w}(i_w)\to B_w,\ O(i_w)\to B_w,\text{ their tower},\\
 \texttt{Oe\_inj},\texttt{Oe\_fin},\texttt{Oe\_loc}
     & \text{injectivity, finiteness and locality from }\rho_w,\\
 \texttt{embW},\texttt{embW\_w}
     & e_w\text{ and the valuation-ring equivalence (NC.7)},\\
 \texttt{eMt},\texttt{eMt\_basis},\texttt{eMt\_lat}
     & \text{the maps and the two conclusions of NC8}.
\end{array}
\tag{NC.20}
\]
The algebra \(O(i_w)\to B_w\) is the actual composite through
\(O_{a_w}(i_w)\), so the scalar-tower field holds as an equality
of those maps. NC4 supplies \texttt{twI\_l}.
For every \(w\in W_{\rm tw}\) and \(1\leq j\leq m\),
\texttt{eMt\_basis} is quantified over every chosen-index basis,
and \texttt{eMt\_lat} is precisely (NC.16) for
\(\mu=\mu_{\rm Tw}(w)\). All 21 fields are therefore proved.

Call the result \(C=\texttt{choiceOfPack}(p,b,\rho)\).
The packet extracted from a choice applies the side convention
once more to \(\texttt{psi'}\). Equation (NC.17) gives the
literal equality
\[
 \texttt{NChoicePack}(C,w)=p_w.
 \tag{NC.21}
\]
Thus the packet, including its selected primes and character,
is recovered exactly; it is not recovered only up to inversion
or isomorphism.

\paragraph{NC10: descending finite-order characters and conjugating a reading.}
Let \(e:L\hookrightarrow K'\) be a field embedding.
If \(L\) contains a primitive \(n\)-th root \(\zeta\), \(n>0\),
and \(\chi:\Gamma\to(K')^\times\) satisfies \(\chi(g)^n=1\),
then there is a unique character \(\psi:\Gamma\to L^\times\)
with \(e(\psi(g))=\chi(g)\).
The \(n\) distinct powers of \(e(\zeta)\) are all the roots of
\(X^n-1\) in \(K'\): a polynomial of degree \(n\) has at most
\(n\) roots. Thus each \(\chi(g)\) has a unique inverse image
among the powers of \(\zeta\).
Injectivity of \(e\) proves that these inverse images preserve
products and the identity.

For a packet \(p=(i,Q,a,\chi)\), the condition
\(\texttt{ValuesIn}(p)\) is exactly the existence of such an
\(n\) and primitive root in \(L\), with \(\chi(\delta)^n=1\)
for every \(\delta\in\mathrm{RCG}(i,Q)\).
Given a reading \(\rho\), apply the preceding construction
to
\(\chi^{[b]}\circ\operatorname{art}_{i,Q}\), followed into
\(\operatorname{Frac}(\rho.O_e)^\times\).
This gives \(\psi_0:G_E\to L^\times\) satisfying
\[
 e_\rho(\psi_0(g))
 =\chi^{[b]}(\operatorname{art}_{i,Q}(g))
 \quad\text{in }\operatorname{Frac}(\rho.O_e).
 \tag{NC.22}
\]
Here and below \(G_E\) is the exact Galois group \(G(E,D.nd.Om)\).
Equation (NC.22) is the literal condition
\(\rho.\texttt{Pulls}(\chi^{[b]},\psi_0)\).

To transport a reading across conjugate places, additionally
assume \(L_0/\mathbb Q\) is Galois. Complex conjugation then
preserves \(L_0\) and \(\mathcal O_{L_0}\).
The supplied two coefficient embeddings satisfy
\[
 s_{i,!b}(x)=s_{i,b}(\bar x),\qquad
 v_{\bar w}(x)=v_w(\bar x).
 \tag{NC.23}
\]
Keep the same DVR and put \(e_{\rho^c}=e_\rho\circ\bar{\phantom{x}}\).
Equations (NC.7) and (NC.23) prove every field of
\(\rho^c:\texttt{Reading}(i,a,!b,\bar w)\).
For a root of unity \(u\in L\), its specified complex image
has absolute value one, so \(\bar u=u^{-1}\).
If (NC.22) holds, the character \(\psi_0\) has finite order
pointwise: the target comes from the finite group
\(\mathrm{RCG}(i,Q)\), and \(e_\rho\) is injective.
Therefore
\[
 e_{\rho^c}(\psi_0(g))
 =e_\rho(\psi_0(g)^{-1})
 =\bigl(\chi^{[b]}(\operatorname{art}(g))\bigr)^{-1}
 =\chi^{[!b]}(\operatorname{art}(g)).
 \tag{NC.24}
\]
The conjugate reading pulls the \emph{same} \(L\)-valued
character, with the side reversed.

\paragraph{NC11: a coherent common character exists under the exact value hypotheses.}
Suppose packets \(t_w\) and sides \(b_w\) are supplied for
\(w\in W_{\rm tw}\), and satisfy the reading and conjugation
parts of \texttt{GoodData}:
\[
 \texttt{ReadsAt}(t_w.i,b_w,w),\quad
 t_{\bar w}=t_w,\quad
 \bar w\ne w\Longrightarrow b_{\bar w}=!b_w.
 \tag{NC.25}
\]
Assume \(L_0/\mathbb Q\) is Galois and
\(\texttt{ValuesIn}(t_w)\) for every such \(w\).
Then there are characters \(\psi_{0,w}:G_E\to L^\times\)
such that
\[
 \psi_{0,\bar w}=\psi_{0,w},\qquad
 w\notin W_{\rm tw}\Longrightarrow\psi_{0,w}=1,
 \tag{NC.26}
\]
and each \(w\in W_{\rm tw}\) has a reading pulling
\(\psi_{0,w}\) from the prescribed packet and side.
This is the exact \texttt{IsCharOf}(t,b,\(\psi_0\)) assertion.

Choose one representative of each orbit of the involution
\(w\mapsto\bar w\). On a twisted representative use NC4
and NC10 to choose a reading and descended character.
On the other member of a two-element orbit use (NC.24);
(NC.25) gives exactly its packet and side.
On a one-element orbit the original choice already suffices:
there is no side-reversal requirement there.
Put the character equal to one on the remaining orbits.
This proves all of (NC.26) and all the required readings.
The additional Hecke-goodness clauses of \texttt{GoodData}
are not used in this character construction.

The result chooses a common character; it does not identify it
with an independently prescribed local Galois character.
That exact compatibility is the next hypothesis,
\texttt{LocsCarry}, and is kept on these same objects.

\paragraph{NC12: the actual tied choice, with all five \texttt{NChTie} fields.}
Retain \(\texttt{ReadElsewhere}(D)\), meaning that every
\(w\notin W_{\rm tw}\) has some \(\texttt{ReadsAt}(i,b,w)\).
On \(W_{\rm tw}\) suppose packets \(t_w\), sides \(b_w\),
the readings in (NC.25), and \(\texttt{RankTie}(t_w,b_w,w)\)
are supplied. Suppose also that a \emph{fixed} family
\(\psi_{0,w}\) satisfies \texttt{IsCharOf}. For each existing
local object \(a:D.nd.Loc\), put \(w_a=D.nd.w(a)\);
the record supplies \(w_a\in W_{\rm tw}\).
The condition \texttt{LocsCarry} requires
\[
 \begin{gathered}
 D.ntw.twOf(a)=t_{w_a}.i,\qquad
 D.ntw.twSide(a)=b_{w_a},\qquad
 D.nd.Q(a)=t_{w_a}.Q,\\
 D.nd.embK(a)(\psi_{0,w_a}(g))
   =D.ntw.psiL(a)(g)
       \quad\text{in }\operatorname{Frac}(D.nd.O(a))
       \quad(g\in G_E).
 \end{gathered}
 \tag{NC.27}
\]
The last line includes the specified injection of the unit
on the right into that fraction field.

For \(w\in W_{\rm tw}\), use the reading supplied by
\texttt{IsCharOf}, so that it pulls this very \(\psi_{0,w}\).
For \(w\notin W_{\rm tw}\), choose a reading index and side
from \texttt{ReadElsewhere}, set
\[
 p_w=(i_w,\varnothing,0,1),
 \tag{NC.28}
\]
and use NC4 with \(k=0\). It pulls \(\psi_{0,w}=1\).
Apply NC9 to this family. Rank equality is needed only
on \(W_{\rm tw}\), where it is the prescribed \texttt{RankTie}.
The resulting \(C:\texttt{NChoice}(D)\) has
\[
 \texttt{NChoicePack}(C,w)=t_w,\quad C.sd(w)=b_w
 \qquad(w\in W_{\rm tw}).
 \tag{NC.29}
\]
This is exactly \texttt{Records}(C,t,b).

For \(\texttt{NChTie}(D,C)\), its first three fields
\texttt{twOf\_eq}, \texttt{twSide\_eq}, and \texttt{Q\_eq}
follow from (NC.27) and (NC.29).
Its fourth field \texttt{psi\_eq} uses the witness
\(\psi_0\) already fixed above: conjugation invariance is
(NC.26), the chosen readings prove its pullback equality
at every \(w\), and the last line of (NC.27) gives its equality
at every existing local object.
Its fifth field \texttt{embedding\_baseField\_condition}
is the denominator equality in (NC.7), restricted to \(z\ne0\).
The composite \(A_0\to O(i_w)\to O_{a_w}(i_w)\to B_w\)
is the same one used in (NC.20). This proves
\[
 \exists\, C:\texttt{NChoice}(D),\qquad
 \texttt{NChTie}(D,C)\ \wedge\ \texttt{Records}(C,t,b).
 \tag{NC.30}
\]
No \texttt{formsIso} or reading-extension existence is
left as an extra premise in (NC.30).

If the full \texttt{GoodData}(I,t,b) is supplied, (NC.29)
also gives \texttt{TwConjAt}(D,C) and
\texttt{TwGoodAt}(D,C,I). For the former, substitute the
identical packets at \(w,\bar w\) and their opposite sides
when distinct. For the latter, substitute the identical
packet in the \texttt{GoodPack} assertion; (NC.17) restores
its normalized character. This transports supplied
Hecke-goodness and does not establish the analytic
\texttt{GoodTw} supplier.

\paragraph{NC13: the exact rank and reading requirements are necessary as well.}
For every supplied \(C:\texttt{NChoice}(D)\) and
\(w\in W_{\rm tw}\),
\[
 \texttt{RankTie}(\texttt{NChoicePack}(C,w),C.sd(w),w)
 \tag{NC.31}
\]
holds. Choose an \(L\)-basis of the relevant \(M_{t,L}\).
The field \texttt{eMt\_basis} makes its images a basis after
extension to \(\operatorname{Frac}(C.Oe(w))\), so the two
dimensions there equal the cardinality of that basis.
NC6 identifies the target dimension with the original
local fraction-field dimension; (NC.17) removes the two
side inversions. This proves (NC.31).

If \(C\) is tied, it also gives
\(\texttt{ReadsAt}(C.twI(w),C.sd(w),w)\) for every \(w\).
For \(x\in\mathcal O_{L_0}\), let \(y\in A_0\) be its
integral image through the side embedding, as in NC2.
The tied denominator identity with denominator \(1\)
identifies \(C.embW(w)(x)\) with the image of \(y\) in \(B_w\).
The two coefficient maps are local, so this image is a
nonunit exactly when \(y\) is a nonunit.
For \(x\ne0\), \texttt{embW\_w} applied to \(x^{-1}\)
identifies nonunit status with \(v_w(x)>0\);
for \(x=0\) both statements hold directly.
Thus the precise maximal-ideal reading (NC.6) follows.
In particular \texttt{ReadElsewhere} follows from a tied
choice. Two tied choices have identical \(i,b,Q\) at each
existing local object, since each equals the corresponding
entry in (NC.27).

These converse facts explain the retained hypotheses.
They cannot be used to prove \texttt{RankTie} before the
choice exists, nor do they assert uniqueness of its DVR,
embedding, character away from the constrained objects,
or chosen lattice isomorphisms.

\paragraph{NC14: finite-order tame characters and nontrivial distinguished twists.}
Now fix the exact \texttt{Stages} \(X\), character frame \(f\),
and \(\texttt{TwDat}\ T\) used in the tuple endpoint.
Write \(E=X.B.host.E\), \(F=X.B.host.F\), and
\[
 \xi_{\rm Tw}(w)
  =\delta_{\operatorname{resChar}(w)}\,T.\psi(w),\qquad
 \delta_\ell=
 \begin{cases}f.\delta'&\ell=X.B.host.pde,\\
               f.\delta&\text{otherwise}.
 \end{cases}
 \tag{NC.32}
\]
Let
\(\Xi_T=\{1,f.\delta,f.\delta'\}\cup
              \{\xi_{\rm Tw}(w):w\}\).
Every member has finite order, is trivial on \(F\) in the
specified finite and infinite restriction sense, and is
tame at every finite place.
Indeed these assertions hold for both distinguished
characters by their \texttt{IsDelta} fields, and for
\(T.\psi(w)\) by its \texttt{IsTwist} field.
For a product, exponents \(n,n'>0\) can be replaced by
\(nn'\); restrictions and local values multiply.
This proves the assertions for (NC.32), as well as for
the identity character.

Every \(\xi_{\rm Tw}(w)\) is nontrivial.
For either distinguished character supported above its
prime \(p\), choose a prime \(P\) of \(E\) above \(p\).
Such a prime exists: \(\mathcal O_E/p\mathcal O_E\) is a
nonzero finite ring, so it has a maximal ideal.
It is nonzero because \(1/p\) is not an algebraic integer.
The twist character \(T.\psi(w)\) is unramified at \(P\).
Otherwise its \texttt{IsTwist} ramification clause says
that \(N(P)\) is a rational prime outside \texttt{offSet}.
As \(P\mid p\), that prime is \(p\), but both distinguished
primes belong to \texttt{offSet}, a contradiction.
The \texttt{IsDelta.ine} clause says that the restriction
of the distinguished character to local units has order
exactly \(15\): exponent \(n\) kills it if and only if
\(15\mid n\). In particular it is nontrivial.
An unramified twist is trivial on those units, so their
product cannot be the identity.

There are only finitely many places of \(L\) above a finite
set of rational primes: each is an ideal of the finite ring
\(\mathcal O_L/p\mathcal O_L\).
Consequently
\[
 \{w:\operatorname{resChar}(w)\in T.pl\}
 \quad\text{and}\quad \Xi_T
 \quad\text{are finite}.
 \tag{NC.33}
\]
For the second assertion, outside the first set
\(T.\psi(w)=1\), so (NC.32) equals one of the two
distinguished characters. Conjugation preserves residue
characteristic, and \(T.\psi(\bar w)=T.\psi(w)\), so
\(\xi_{\rm Tw}(\bar w)=\xi_{\rm Tw}(w)\) as well.

\paragraph{NC15: \texttt{MuIsW} identifies the exact finite support.}
Retain the supplied \(\texttt{MuIsW}(X,f,T,W,\xi)\).
Its ten clauses are: injectivity and range \(\Xi_T\) of
\(\xi\); its specified values \(1,f.\delta,f.\delta'\)
at \(\mu_0,\mu_\star,\mu_{\star'}\); the identity
\(\xi(\mu_{\rm Tw}(w))=\xi_{\rm Tw}(w)\); the assignments
\(\mu_{\rm At}(w)=\mu_{\rm Tw}(w)\) on \(W\) and
\(\mu_{\rm At}(w)=\mu_0\) off \(W\); the inclusion
of \(W\) into the first set of (NC.33); and
\(T.\psi(w)=1\) off \(W\).
Then
\[
 \mu_{\rm At}(w)\ne\mu_0\quad\Longleftrightarrow\quad w\in W,
 \qquad W\text{ finite},\qquad
 w\in W\Longrightarrow\bar w\in W.
 \tag{NC.34}
\]
If \(w\notin W\), its off-support assignment gives
\(\mu_{\rm At}(w)=\mu_0\).
If \(w\in W\) and this equality held, applying \(\xi\)
would give \(\xi_{\rm Tw}(w)=1\), contrary to NC14.
This proves the equivalence. Finiteness follows from
the supplied inclusion and (NC.33).
The row field \(\mu_{\rm At}(\bar w)=\mu_{\rm At}(w)\)
then proves conjugation stability.
In particular \(W\) is the exact support, although it
need not be the set of all places above \(T.pl\).

\paragraph{NC16: all twelve fields of \texttt{TupTied} on the same characters.}
Under the hypotheses of NC15, let
\(C:\texttt{NChoice}(X.toD7)\) satisfy
\[
 C.Q(w)=T.Q(w)\qquad(w\in W).
 \tag{NC.35}
\]
There is a
\(\mathrm{tp}:\texttt{TupTied}(X.toD7,\texttt{Fr.at}(X,f),C)\)
with the literal equalities
\[
 \mathrm{tp}.\xi=\xi,\qquad
 \mathrm{tp}.\psi=T.\psi.
 \tag{NC.36}
\]
Its twelve fields are proved as follows.
The fields \texttt{xi} and \texttt{psi} are those exact
functions. \texttt{xi\_inj} is the \texttt{MuIsW} injection.
For \texttt{xi\_fin}, \texttt{xi\_triv}, and
\texttt{xi\_tame}, the range identity puts every
\(\xi(\mu)\) in \(\Xi_T\), and NC14 gives the three assertions.
The field
\texttt{trueColumnChoices\_heckeChar\_reference\_eq\_one}
and the two fields \texttt{xi\_star}, \texttt{xi\_star'}
are exactly the three specified values in NC15.
The field \texttt{xi\_tw} is (NC.32) and the supplied
twisted-tuple identity, since \texttt{Fr.at} uses the same
distinguished characters.
For \texttt{psi\_is}, a nonreference place belongs to \(W\)
by (NC.34); substitute (NC.35) in the actual
\(T.\texttt{psi\_is}(w)\) assertion.
For \texttt{psi\_un}, a reference place is outside \(W\),
and the support clause of \texttt{MuIsW} gives \(T.\psi(w)=1\).
This proves all twelve fields and both equalities (NC.36).

No equality of the coefficient indices \(C.k'(w)\) and
\(T.k(w)\) is required by this endpoint. Nor does it identify
the Galois character \(\psi_{0,w}\) of NC12 with the Hecke
character \(T.\psi(w)\) without a supplied reciprocity
compatibility. They are different types of data.

\paragraph{NC17: the same-object join and the precise remaining supplier boundary.}
Combine NC12 with NC15--NC16 for the same \(D=X.toD7\).
Suppose the twisted packet has
\[
 t_w.Q=T.Q(w)\qquad(w\in W=W_{\rm tw}).
 \tag{NC.37}
\]
Taking the \(Q\)-component of (NC.29) proves (NC.35).
Thus NC12 and NC16 construct simultaneously
\[
 \begin{gathered}
 C:\texttt{NChoice}(D),\qquad \texttt{NChTie}(D,C),\\
 \mathrm{tp}:\texttt{TupTied}(D,\texttt{Fr.at}(X,f),C),
 \qquad \mathrm{tp}.\xi=\xi,\quad \mathrm{tp}.\psi=T.\psi.
 \end{gathered}
 \tag{NC.38}
\]
The component entries \(C.Q,C.k',C.sd,C.twI\) on \(W\)
are the ones in \(t,b\). Off \(W\), the empty-prime,
exponent-zero, trivial-character packet (NC.28) supplies
the unconstrained entries.

For an instance of the previously reconstructed TX16 endpoint,
also retain the literal frame equality
\(I.\texttt{cf}=\texttt{Fr.at}(X,f)\) and
\(\xi=I.\texttt{xi}\), together with its already stated
geometric and stage premises. Equation (NC.36) then gives
\(\mathrm{tp}.\xi=I.\texttt{xi}\) on the same objects,
so (NC.38) supplies its \texttt{NChoice}/\texttt{TupTied}
arguments. TX16's geometric conclusion is therefore
available under these more explicit arithmetic premises.
No new comparison of independently chosen frames or
character families is silently inserted.

The complete conditional suppliers proved here are
\texttt{exists\_coeff\_ext},
\texttt{ReadsAt.nonempty\_reading},
\texttt{exists\_formsIso},
\texttt{choiceOf}/\texttt{choiceOfPack},
\texttt{exists\_choice},
\texttt{exists\_tied\_choice\_of\_reads}, and
\texttt{tupTied\_of\_muIsW}, including their finite-support
and normalization consequences. For \texttt{exists\_choice}
alone, \texttt{IsCharOf} and \texttt{LocsCarry} are unnecessary:
use arbitrary readings from NC4 before NC9.

The remaining simultaneous supply is explicit:
the actual \texttt{D7} row and local models;
twisted packets and sides which satisfy \texttt{ReadsAt}
and \texttt{RankTie}; \texttt{ReadElsewhere} or its
\texttt{TwCover} supplier; a common family satisfying
\texttt{IsCharOf} and the same-object \texttt{LocsCarry};
and \texttt{MuIsW} with the exact prime equality (NC.37)
and, where TX16 is used, the stated frame and \(\xi\)
identifications.
NC11 can supply \texttt{IsCharOf} under its primitive-root,
conjugation, and \(L_0/\mathbb Q\)-Galois hypotheses,
but it does not supply \texttt{LocsCarry} for that new family.
All constructions are relative to these data; this section
does not prove a full \texttt{Later}, \texttt{Pins}, host
window, or final RH/GRH existence theorem.
